% Historical source (removed by author request): references/_OLD_Übungsaufgaben/Übungsblatt 5 - Differenzialrechnung/Blatt5.tex:258
\begin{ExerciseBlock}
\ExerciseHeading{12}{Optimierung: Gassen Speyers}{differential_calculus_alley_transport}
Ein nicht teilbarer und aufstellbarer länglicher Gegenstand soll durch die Gassen Speyers zu einem Technik-Museum transportiert werden. Die engste Stelle auf dem Weg wird durch zwei zueinander rechtwinklig verlaufende Gassen der Breite $a$ und $b$ gebildet. Berechnen Sie die maximale Länge des Gegenstandes, wenn $a=8\,$m und $b=4\,$m ist.\par
\begin{center}
	\CachedTikzFigure{content/resources/chapters/differential_calculus/tikz/alley_transport.tex}
\end{center}

\end{ExerciseBlock}
\UseExerciseSolution{differential_calculus_alley_transport}
